\chapter{Experiments and Results}
\label{ch:experiments-results}

\section{Introduction}

The previous chapter described the design of our \gls{rl}-based auto-scheduler across four incremental versions, each introducing a specific improvement over the baseline system: hardware-aware observation (V1), shaped reward (V2), transformer-based loop nest encoder (V3), and the unified combination with expanded action space (V4). This chapter evaluates these versions experimentally.

The evaluation is designed around a set of controlled comparisons. For each novel feature, the corresponding version is compared directly against the baseline (V0) on identical benchmarks, using identical training protocols and hyperparameter configurations. This isolates the contribution of each individual design choice. The unified version (V4) is then compared against the baseline and against each individual version to measure the cumulative benefit of combining all improvements.

The chapter is organized as follows. Section~5.2 describes the experimental setup: hardware platform, software stack, training protocol, and dataset composition. Section~5.3 defines the evaluation metrics. Sections~5.4 through~5.8 present the results, starting with baseline performance and proceeding through each version. Section~5.9 presents results on full neural network models. Section~5.10 provides a discussion of the findings, and Section~5.11 concludes.

\section{Experimental Setup}

\subsection{Hardware and Software Environment}

All experiments are conducted on a Slurm-based HPC cluster. The primary compute nodes are equipped with dual-socket Intel Xeon processors at 2.4~GHz, with 32~physical cores per socket, 256~GB of RAM, and three levels of cache (L1: 32~KB per core, L2: 1~MB per core, L3: 33~MB shared across all cores). Vectorization is supported through the AVX-512 instruction set. Some experiments are repeated on an AMD EPYC processor to evaluate cross-platform generalization.

The software stack consists of Python~3.11, PyTorch~2.0~\cite{paszke2019pytorch}, MLIR built from the LLVM~19.x release branch, and the \gls{mlir} Python bindings. All \gls{mlir} benchmarks are compiled through the progressive lowering pipeline (Linalg to Affine to \gls{scf} to Vector to \gls{llvm}) and \gls{jit}-compiled at optimization level \texttt{-O3} using the MLIR ExecutionEngine. Execution time measurements are performed on zero-initialized buffers to eliminate data-dependent variability, with a warm-up invocation preceding each timed invocation.

\subsection{Training Protocol and Hyperparameters}

All versions share a common training protocol based on the \gls{ppo} algorithm~\cite{schulman2017ppo} with Generalized Advantage Estimation~\cite{schulman2015gae}. Training proceeds in iterations, each consisting of trajectory collection under the current policy followed by a \gls{ppo} update using the Adam optimizer~\cite{kingma2014adam}.

\begin{table}[H]
    \centering
    \caption{Training hyperparameters shared across all versions.}
    \label{tab:hyperparameters}
    \begin{tabularx}{\textwidth}{l l l}
        \toprule
        \textbf{Parameter} & \textbf{Value} & \textbf{Description} \\
        \midrule
        Learning rate & $3 \times 10^{-4}$ & Adam optimizer step size \\
        \gls{ppo} clip range ($\epsilon$) & 0.2 & Policy ratio clipping bound \\
        Discount factor ($\gamma$) & 1.0 & GAE discount (no time preference) \\
        GAE~$\lambda$ & 0.95 & Bias-variance trade-off \\
        \gls{ppo} epochs per iteration & 4 & Number of update epochs on collected data \\
        \gls{ppo} mini-batch size & 64 & Batch size for gradient estimation \\
        Value loss coefficient ($c_1$) & 0.5 & Weight of value loss in combined objective \\
        Entropy coefficient ($c_2$) & 0.01 & Exploration bonus weight \\
        Max gradient norm & 0.5 & Gradient clipping threshold \\
        Training iterations & 2000 & Total number of trajectory-collection cycles \\
        Benchmarks per trajectory & 32 & Number of programs sampled per iteration \\
        Truncation length & 5 & Maximum transformations per episode \\
        \bottomrule
    \end{tabularx}
\end{table}

Table~\ref{tab:hyperparameters} lists the hyperparameters shared across all versions. The discount factor $\gamma = 1.0$ reflects the episodic nature of compiler optimization: the agent's objective is to maximize total speedup rather than to minimize the number of steps. The entropy coefficient of $0.01$ provides a moderate exploration incentive. Gradient clipping at $0.5$ prevents destabilizing updates.

Version-specific hyperparameters are detailed where relevant. V3 and V4, which use a transformer encoder, employ a model with $d_\text{model} = 256$, $h = 8$ attention heads, $N = 3$ encoder layers, feedforward dimension $d_\text{ff} = 1024$, dropout of $0.1$, and GELU activation. V4 additionally enables the expanded action space with three tile sizes per dimension, three pad multiples (powers of two), and three unroll factors (powers of two). The shaped reward in V2 and V4 uses weights of $1.0$ for arithmetic intensity, $0.1$ for vectorizability, and $0.1$ for parallel-loop ratio, with a global shaping scale of $1.0$ and clipping at $\pm 2.0$.

\section{Evaluation Metrics}

The primary evaluation metric is the execution-time speedup, defined as the ratio of the unoptimized baseline execution time to the optimized execution time:

\begin{equation}
    \text{Speedup} = \frac{T_\text{baseline}}{T_\text{optimized}}
\end{equation}

where $T_\text{baseline}$ is the execution time of the unoptimized \gls{mlir} code and $T_\text{optimized}$ is the execution time after applying the transformation schedule selected by the agent. A speedup greater than $1.0$ indicates an improvement over the baseline; a speedup below $1.0$ indicates a degradation.

For aggregated comparisons across multiple benchmarks, we report the geometric mean of the per-benchmark speedups:

\begin{equation}
    \text{Geometric Mean Speedup} = \left( \prod_{i=1}^{n} \text{Speedup}_i \right)^{1/n}
\end{equation}

The geometric mean is preferred over the arithmetic mean because speedup ratios are multiplicative in nature. A benchmark with a $10\times$ speedup and another with a $0.1\times$ slowdown should average to $1.0\times$, which the geometric mean correctly captures.

Additional metrics include the best achieved speedup (the maximum speedup obtained across multiple evaluation episodes for a single benchmark), the fraction of benchmarks that achieve a speedup greater than $1.0$, and training convergence metrics (cumulative episode reward, policy loss, value loss, and policy entropy as a function of training iteration).

\section{Baseline Performance}

The baseline version (V0) corresponds to the system described by Tirichine et al.~\cite{tirichine2025reinforcement}. It employs an \gls{ast}-based state representation with a recursive \gls{lstm} encoder, a sparse terminal reward (log-speedup at episode completion), and an action space of six transformation types (Tiling, TiledParallelization, TiledFusion, Interchange, Vectorization, NoTransformation). The baseline is trained on the full dataset for 2000 iterations and evaluated on the held-out evaluation set.

\begin{table}[H]
    \centering
    \caption{Baseline (V0) speedups on the evaluation set. Results are aggregated by benchmark category.}
    \label{tab:v0-baseline}
    \begin{tabularx}{\textwidth}{l c c c}
        \toprule
        \textbf{Benchmark Category} & \textbf{Number of Benchmarks} & \textbf{Geometric Mean Speedup} & \textbf{Best Speedup} \\
        \midrule
        Matrix multiplication (various sizes) & 12 & -- & -- \\
        2D convolution (various sizes) & 10 & -- & -- \\
        Element-wise operations & 8 & -- & -- \\
        Operation sequences (2--5 ops) & 15 & -- & -- \\
        \midrule
        \textbf{Overall} & 45 & -- & -- \\
        \bottomrule
    \end{tabularx}
\end{table}

Table~\ref{tab:v0-baseline} reports the baseline speedups across benchmark categories ...

\section{Evaluation: Hardware-Aware Observation}

Version V1 introduces a hardware-aware observation space. The agent's input is augmented with a fixed-size vector encoding the target processor's cache hierarchy (L1, L2, L3 cache sizes in kilobytes), core count (physical and logical), \gls{simd} width in bits, and clock frequency in megahertz. These values are either auto-detected from the compute node or manually specified via configuration. The policy architecture, action space, and reward function are otherwise identical to V0.

\subsection{Experiment Design}

V1 is trained and evaluated under the same protocol as V0, on the same compute node whose hardware characteristics are encoded in the observation. The comparison isolates the effect of hardware awareness: any difference in performance between V0 and V1 can be attributed to the agent's ability to condition its decisions on known platform parameters.

To evaluate cross-platform generalization, both V0 and V1 are evaluated on a different hardware configuration (an AMD EPYC node with distinct cache sizes and core count) without retraining. A policy that learns platform-appropriate strategies implicitly (V0) is expected to transfer poorly, while a policy that conditions on explicit hardware features (V1) is expected to retain more of its performance.

\begin{table}[H]
    \centering
    \caption{Speedup comparison: V0 versus V1 on the primary training hardware.}
    \label{tab:v1-comparison}
    \begin{tabularx}{\textwidth}{l c c c}
        \toprule
        \textbf{Benchmark Category} & \textbf{V0 Geometric Mean} & \textbf{V1 Geometric Mean} & \textbf{Relative Change (\%)} \\
        \midrule
        Matrix multiplication & -- & -- & -- \\
        2D convolution & -- & -- & -- \\
        Element-wise operations & -- & -- & -- \\
        Operation sequences & -- & -- & -- \\
        \midrule
        \textbf{Overall} & -- & -- & -- \\
        \bottomrule
    \end{tabularx}
\end{table}

\begin{table}[H]
    \centering
    \caption{Cross-platform generalization: V0 and V1 evaluated on a different hardware platform without retraining.}
    \label{tab:v1-cross-platform}
    \begin{tabularx}{\textwidth}{l c c}
        \toprule
        \textbf{Version} & \textbf{Geometric Mean Speedup (Training HW)} & \textbf{Geometric Mean Speedup (Different HW)} \\
        \midrule
        V0 (hardware-agnostic) & -- & -- \\
        V1 (hardware-aware) & -- & -- \\
        \bottomrule
    \end{tabularx}
\end{table}

\subsection{Analysis}

...

\section{Evaluation: Shaped Reward Function}

Version V2 replaces the sparse terminal reward of V0 with a shaped reward function that provides dense intermediate feedback at every transformation step. The shaped reward consists of two components: the terminal speedup reward applied at episode completion (identical to V0), and a per-step shaped delta computed from static code analysis of the \gls{mlir} code before and after each transformation. The shaped delta is based on three static metrics:

\begin{itemize}
    \item \textbf{Arithmetic intensity}: the ratio of floating-point operations to bytes transferred, estimated from the operation's loop structure and access patterns. Higher arithmetic intensity indicates better compute utilization.
    \item \textbf{Vectorizability}: a heuristic score reflecting the fraction of loop iterations that can be mapped to \gls{simd} vector lanes given the current access strides and loop bounds.
    \item \textbf{Parallel-loop ratio}: the fraction of loops that carry no data dependencies and can therefore be parallelized.
\end{itemize}

The per-step reward at time $t$ is:

\begin{equation}
    r_t^\text{shaped} = w_\text{ai} \cdot \Delta\text{AI} + w_\text{vec} \cdot \Delta\text{Vec} + w_\text{par} \cdot \Delta\text{Par}
\end{equation}

where $\Delta\text{AI}$, $\Delta\text{Vec}$, and $\Delta\text{Par}$ are the changes in the respective static metrics induced by the transformation at step $t$, and the weights are $w_\text{ai} = 1.0$, $w_\text{vec} = 0.1$, $w_\text{par} = 0.1$. The shaped component is clipped to $[-2.0, 2.0]$ and added to the terminal speedup reward at the end of the episode.

\subsection{Experiment Design}

V2 is trained using the shaped reward and compared against V0 (sparse terminal reward only). Both versions share the same state representation, action space, policy architecture, and hyperparameters, isolating the effect of the reward function. In addition to per-benchmark speedups, the training dynamics are compared: the number of iterations required to reach a given performance threshold, the stability of the learning curve, and the variance of the achieved speedup across multiple evaluation seeds.

\begin{table}[H]
    \centering
    \caption{Speedup comparison: V0 (sparse reward) versus V2 (shaped reward).}
    \label{tab:v2-comparison}
    \begin{tabularx}{\textwidth}{l c c c}
        \toprule
        \textbf{Benchmark Category} & \textbf{V0 Geometric Mean} & \textbf{V2 Geometric Mean} & \textbf{Relative Change (\%)} \\
        \midrule
        Matrix multiplication & -- & -- & -- \\
        2D convolution & -- & -- & -- \\
        Element-wise operations & -- & -- & -- \\
        Operation sequences & -- & -- & -- \\
        \midrule
        \textbf{Overall} & -- & -- & -- \\
        \bottomrule
    \end{tabularx}
\end{table}

\subsection{Analysis}

...

\section{Evaluation: Transformer Loop-Nest Encoder}

Version V3 replaces the recursive \gls{lstm} encoder of V0 with a transformer-based encoder. Each loop level in the nest is treated as a token, with positional encodings capturing the depth and role (outermost, intermediate, innermost) of each loop in the nest hierarchy. The tokens are processed through $N = 3$ transformer encoder layers with $h = 8$ attention heads and model dimension $d_\text{model} = 256$. The resulting token embeddings are pooled (via a \texttt{[CLS]} token or mean pooling) into a fixed-size vector that feeds into the policy backbone. The state representation, action space, and reward function are otherwise identical to V0.

\subsection{Experiment Design}

V3 is compared against V0 on the same benchmarks and training protocol. The comparison isolates the effect of the encoder architecture. To test the hypothesis that the transformer's self-attention mechanism better captures long-range dependencies in deeply nested loops, the evaluation benchmarks are stratified by loop nest depth. Shallow nests (two to three loops) serve as a control where the sequential LSTM should perform adequately; deep nests (four or more loops) are the test case where the transformer's ability to attend directly between any pair of loop levels is expected to provide an advantage.

\begin{table}[H]
    \centering
    \caption{Speedup comparison: V0 (LSTM) versus V3 (Transformer), stratified by loop nest depth.}
    \label{tab:v3-comparison}
    \begin{tabularx}{\textwidth}{l c c c}
        \toprule
        \textbf{Nest Depth} & \textbf{V0 Geometric Mean} & \textbf{V3 Geometric Mean} & \textbf{Relative Change (\%)} \\
        \midrule
        Shallow (2--3 loops) & -- & -- & -- \\
        Deep (4+ loops) & -- & -- & -- \\
        \midrule
        \textbf{Overall} & -- & -- & -- \\
        \bottomrule
    \end{tabularx}
\end{table}

\subsection{Analysis}

...

\section{Evaluation: Expanded Action Space and Unified System}

Version V4 combines the three preceding improvements (hardware-aware observation, shaped reward, transformer encoder) and further introduces three additional transformation types: Pad, Pack, and Unroll. The Pad transformation adjusts tensor dimensions by adding padding elements to align memory accesses with \gls{simd} vector widths. The Pack transformation reorganizes data layout for improved cache utilization. The Unroll transformation replicates the loop body to reduce loop overhead and expose instruction-level parallelism, with the unroll factor selected from powers of two.

\subsection{Experiment Design}

V4 is evaluated against V0 to measure the cumulative impact of all four improvements combined. In addition, V4 is compared against each individual version (V1, V2, V3) to assess the marginal benefit of combining features. All versions are trained and evaluated under identical protocols.

\begin{table}[H]
    \centering
    \caption{Cumulative comparison: V0 versus V4 (unified system).}
    \label{tab:v4-comparison}
    \begin{tabularx}{\textwidth}{l c c c}
        \toprule
        \textbf{Benchmark Category} & \textbf{V0 Geometric Mean} & \textbf{V4 Geometric Mean} & \textbf{Relative Change (\%)} \\
        \midrule
        Matrix multiplication & -- & -- & -- \\
        2D convolution & -- & -- & -- \\
        Element-wise operations & -- & -- & -- \\
        Operation sequences & -- & -- & -- \\
        \midrule
        \textbf{Overall} & -- & -- & -- \\
        \bottomrule
    \end{tabularx}
\end{table}

\begin{table}[H]
    \centering
    \caption{Per-version comparison: geometric mean speedups across all evaluation benchmarks.}
    \label{tab:all-versions}
    \begin{tabularx}{\textwidth}{l c c c c c}
        \toprule
        \textbf{Benchmark Category} & \textbf{V0} & \textbf{V1} & \textbf{V2} & \textbf{V3} & \textbf{V4} \\
        \midrule
        Matrix multiplication & -- & -- & -- & -- & -- \\
        2D convolution & -- & -- & -- & -- & -- \\
        Element-wise & -- & -- & -- & -- & -- \\
        Op. sequences & -- & -- & -- & -- & -- \\
        \midrule
        \textbf{Overall} & -- & -- & -- & -- & -- \\
        \bottomrule
    \end{tabularx}
\end{table}

\subsection{Analysis}

...

\section{Full Neural Network Model Results}
...

\begin{table}[H]
    \centering
    \caption{Speedups on full neural network models. Results include standard PyTorch execution and PyTorch JIT as competitive baselines.}
    \label{tab:full-models}
    \begin{tabularx}{\textwidth}{l c c c c c c}
        \toprule
        \textbf{Model} & \textbf{V0} & \textbf{V1} & \textbf{V2} & \textbf{V3} & \textbf{V4} & \textbf{PyTorch} \\
        \midrule
        VGG-16 & -- & -- & -- & -- & -- & -- \\
        ResNet-50 & -- & -- & -- & -- & -- & -- \\
        MobileNetV2 & -- & -- & -- & -- & -- & -- \\
        \bottomrule
    \end{tabularx}
\end{table}

Table~\ref{tab:full-models} presents the speedups achieved on the three full models ...

\section{Discussion}

...

\textbf{Contribution of individual features.} \\
...

\textbf{Learning efficiency.} \\
...

\textbf{Limitations.} \\
...

\textbf{Comparison with prior work.} \\
...

\section{Conclusion}

This chapter has presented a systematic experimental evaluation of the four versions of the \gls{rl}-based \gls{mlir} auto-scheduler introduced in Chapter~4. Each version was compared against the baseline (V0) on a shared set of benchmarks, under identical training protocols, to isolate the contribution of each design choice. The evaluation covered single operations, operation sequences, and full neural network models (gpt-2, ConvNext, VGG, ...).

The hardware-aware observation (V1) was evaluated on its ability to improve per-benchmark speedups and to enable cross-platform policy transfer. The shaped reward (V2) was assessed both for its effect on final optimization quality and for its impact on training convergence speed. The transformer encoder (V3) was analyzed with a focus on deeply nested loops where long-range dependencies are critical. The unified system (V4) combined all improvements and expanded the action space, yielding cumulative benefits quantified through a comprehensive version comparison and contribution ablation.

The results provide evidence for each proposed improvement and for their combined effect, while also revealing limitations and directions for further work. The next chapter concludes the thesis with a summary of contributions and a discussion of future research directions.
